\scanpage{296}
\mspage{212}

$F$ dans $\operatorname{Ind}(C)\subset C^\wedge$, on sait que $C_{/F}$ est filtrante, et il faut voir qu'\unsure{ille|il|elle} est filtrante grande devant $\pi_0$. On est ramené aux trois cas où $J=\varnothing$, $J=e\amalg e$, et $J=$

\diagram{
\begin{tikzcd}[ampersand replacement=\&]
\cdot\arrow[dr]\& \&\cdot\arrow[dl]\\
\&\cdot\&.
\end{tikzcd}
}

Les deux premiers cas sont essentiellement triviaux (car \unsure{$C$}, et un produit de deux cat. gdes devant $\pi_0$, sont gdes devant $\pi_0$), le troisième cas par contre \unsure{n'a|n'est} pas \inserted{l'air} \unsure{triviale|trivial} a priori. C'est lié visiblement à la question des représentations indicielles de foncteurs $J\longrightarrow\operatorname{Ind}_{\pi_0}(C)$, pour $\operatorname{card}\operatorname{red}J\leq\pi_0$. -- i.e. à la validité de 4.4.13 b), en dehors de l'hypothèse que $C$ soit \unsure{stable} par $\varinjlim$ de cardinal $\leq\pi_0$. Il nous faudra y revenir très prochainement, car cette question m'apparait à présent une question-clef, pour les propriétés de stabilité de la notion de catégorie accessible (en dehors d'\inserted{aucune} hypothèse d'existence de $\varinjlim$ ou $\varprojlim$).

\margin{\unsure{Z|2} \uline{Pas} \uline{prouvé} Mais \unsure{OK} si $\pi_0>\operatorname{card}\unsure{\operatorname{red}|\operatorname{sec}} C$ \unsure{$\Rightarrow$} cf. 4.15.10 p.~\unsure{463|462|403}.}% slanted

\uline{Proposition} 4.8.13. \deleted{\unsure{Soient}} Soit $C$ une catégorie équivalente à une petite catégorie, et \deleted{soit $\pi_0$} soit $\deleted{\unsure{C\geq}}\inserted{\pi_0\geq}\operatorname{card}\operatorname{red}C$ \deleted{un cardinal} \inserted{un cardinal} infini. Alors $M=\operatorname{Ind}_{\pi_0}C\deleted{\unsure{=}}=M_{\pi_0}$, \deleted{et c'est} (i.e. tous les objets de $\operatorname{Ind}_{\pi_0}C$ sont $\pi_0$-accessibles) et $M$ est une enveloppe de Karoubi de $\unsure{\widetilde C|\overline C|C}$,
% archive sheet number: 295
